\chapter{有限回放条件下的医学影像持续学习研究}
\label{chap:task-level-continual-learning}
\setlength{\emergencystretch}{1em}

\section{引言}
\label{sec:task-cl-introduction}

第三章表明，模型在新任务上取得较好结果、保留旧任务能力和泛化到未见任务，是需要分别检验的三个目标；有限回放能够改善其中一些目标，但效果受到回放容量和任务差异影响。本章据此研究有限历史访问下的医学图像分割与配准：在不能重新使用全部历史训练集的条件下，如何利用少量历史样本协调当前任务学习与旧任务保持。分割部分进一步限制当前监督为涂鸦，配准部分则把研究对象扩展为不同器官和模态之间的空间对应关系。

具体而言，本章关注两类输出形式不同、但均受到训练信息访问条件约束的医学影像任务。第一类是 ScribbleCL。本文沿用第三章的 Domain-CL、Class-CL 和 Organ-CL，不重新定义数据集、任务顺序、病例级划分、预处理、输出组织和评价指标。相较于第三章的密集监督基准，当前训练阶段只提供由同一训练标签生成的固定涂鸦；所提 ZSDERpp 进一步允许从容量受限的回放集合读取少量历史图像--涂鸦对和参考特征，并在 Class-CL 中校正背景语义\cite{zhang2026zscribbleseg,cermelli2020mib}。由此，第三章提供统一的场景与评价基础，第四章考察相同任务序列在监督信息进一步受限时的持续学习表现。本文的定量分析集中于 Domain-CL。


第二类研究是持续医学图像配准。与分割预测离散语义标签不同，配准模型需要从固定图像和移动图像估计连续空间变换。当脑部 MRI、腹部 CT、肺部 CT 和腹部 MR--CT 等任务依次到达时，图像外观、解剖区域、模态关系和损失定义均可能发生变化。SAMCL 采用有限经验回放保存少量历史图像对，并通过元持续学习协调当前样本与历史样本的更新方向；在此基础上，锐度感知优化进一步约束参数邻域中的最坏损失，以改善顺序训练中的前向泛化\cite{wang2024samcl,foret2021sam}。因此，SAMCL 采用有限回放的集中式持续学习设定。

将两项研究放在同一章，是因为它们共同回答“有限历史样本应如何参与持续更新”这一问题。ScribbleCL 从监督内容出发，在回放样本上同时保持涂鸦监督和历史特征；SAMCL 从优化过程出发，协调当前与历史图像对的更新方向，并考察前向泛化。两者采用不同网络、损失和评价单位，不构成分割结果输入配准网络的串联流程，也没有在本章中进行联合训练。本章先分别建立证据，再比较二者对有限回放的使用方式；下一章进一步讨论不允许回放历史原始样本且需要跨中心协作时的知识保护。


\section{问题定义与研究动机}
\label{sec:task-cl-formulation}

本章将共同约束与任务特定设计分开描述。共同约束是数据按任务顺序到达、模型持续更新、历史原始样本只能保留有限集合；共同评价问题是当前任务是否学好、旧任务是否退化，以及未训练任务上的直接表现如何。任务特定部分则包括分割的像素语义与涂鸦标签，以及配准的图像对应关系、解剖标签或地标监督。使用同一种回放原则并不意味着两类样本可以混合为一个回放集合，或两类损失可以直接相加。

\subsection{涂鸦监督分割及其持续扩展}
\label{subsec:task-cl-scribble-formulation}

沿用第\ref{chap:foundations}章的分割记号，设医学图像为
$\mathbf{I}:\Omega\rightarrow\mathbb{R}^{C}$，分割模型
$f_{\theta}$ 在位置 $\mathbf{u}\in\Omega$ 输出类别概率
$\mathbf{p}_{\theta}(\mathbf{u})$。在涂鸦监督条件下，只有位置集合
$\Omega_L\subset\Omega$ 具有类别标签，未标注集合为
$\Omega_U=\Omega\setminus\Omega_L$。式\eqref{eq:foundations-partial-ce}
给出的部分交叉熵只在 $\Omega_L$ 上产生直接监督，避免把未知像素简单视为背景。然而，当 $|\Omega_L|\ll|\Omega|$ 时，已标注位置不能直接给出目标的完整区域、轮廓和连通关系。模型即使在涂鸦位置上预测正确，也可能在未标注区域形成欠分割、轮廓偏移或碎片化结果。

静态涂鸦监督学习在固定训练分布上优化单一模型。设训练集为
$\mathcal{D}^{\mathrm{scr}}=\{(\mathbf{I}_i,\mathbf{Y}^{\mathrm{scr}}_i)\}_{i=1}^{N}$，其中
$\mathbf{Y}^{\mathrm{scr}}_i$ 只在 $\Omega_{L,i}$ 上定义。一般目标可以写为
\begin{equation}
\widehat{\theta}
=
\operatorname*{arg\,min}_{\theta}
\frac{1}{N}\sum_{i=1}^{N}
\left[
\mathcal{L}_{\mathrm{PCE}}^{(i)}(\theta)
+
\mathcal{R}_{\mathrm{scr}}^{(i)}(\theta)
\right],
\label{eq:scribble-static-objective}
\end{equation}
其中，$\mathcal{R}_{\mathrm{scr}}$ 表示从输入变换、预测一致性或空间结构中获得的附加约束。该目标的关键不是把未标注像素补写为确定标签，而是既保持涂鸦所提供的监督范围，又为未标注区域引入可解释且受条件限制的训练信号。

在此基础上，ScribbleCL 进一步考虑任务按阶段到达。设任务序列为
$\mathcal{T}=\{1,\ldots,T\}$，第 $t$ 个阶段获得涂鸦训练集
$\mathcal{D}^{\mathrm{scr}}_t$，新增类别集合为 $\mathcal{C}_t$，截至该阶段的已见类别集合为
$\mathcal{C}_{1:t}=\bigcup_{\tau=1}^{t}\mathcal{C}_{\tau}$。完成第 $t$ 个阶段后，模型参数记为 $\theta_t$。令
$d_{t,j}$ 表示模型 $\theta_t$ 在任务 $j$ 测试集上的平均 Dice，则
$\mathbf{D}^{\mathrm{scr}}=[d_{t,j}]$ 沿用第\ref{chap:benchmark}章的阶段--任务性能矩阵。矩阵中的测试任务、密集测试标签和评价时点均与第三章相同，区别在于 $\theta_t$ 是在涂鸦监督和相应历史访问条件下获得的。与密集监督 Class-CL 不同，此时主对角线 $d_{t,t}$ 既受到持续约束，也受到涂鸦监督不足的影响；旧任务条目则共同反映历史保持和早期弱监督模型本身的上限。因此，只报告最终平均 Dice 无法判断性能损失主要来自当前任务没有学好，还是来自后续训练遗忘了早期任务。

在历史访问方面，ScribbleCL 不允许完整访问全部历史训练集，但 ZSDERpp 可以从容量受限的回放集合重新读取少量历史图像--涂鸦对及其参考特征。当前阶段的弱监督目标由第\ref{subsec:prior-scribble-foundation}节给出，历史阶段的损失由同一部分交叉熵、全局一致性和空间先验在回放样本上重新计算，并加入特征保持约束。Domain-CL 使用共享输出头，Organ-CL 依据任务标识选择独立输出头，Class-CL 在统一目标内调整背景概率的语义。完整的阶段--任务性能矩阵用于区分当前监督不足、历史保持与前向泛化。


\subsection{持续医学图像配准}
\label{subsec:task-cl-registration-formulation}

医学图像配准的基本映射与训练目标已在式\eqref{eq:foundations-reg-learning}中给出。设第 $t$ 个配准任务的数据集为
$\mathcal{D}^{(t)}=\{(I_{F,i}^{(t)},I_{M,i}^{(t)})\}_{i=1}^{N_t}$，配准网络
$g_{\theta}$ 输出从固定图像域到移动图像域的变换
$\phi_{\theta}^{(t)}=g_{\theta}(I_F^{(t)},I_M^{(t)})$。不同任务可以具有不同模态、解剖区域和监督信息，因此其配准损失记为
$\mathcal{L}^{(t)}(\phi_{\theta};I_F,I_M)$。若所有已见任务数据可以一并访问，则完成前 $K$ 个任务后的理想累积目标为
\begin{equation}
\widehat{\theta}^{(K)}
=
\operatorname*{arg\,min}_{\theta}
\sum_{t=1}^{K}
\mathbb{E}_{(I_F,I_M)\sim\mathcal{D}^{(t)}}
\left[
\mathcal{L}^{(t)}(\phi_{\theta};I_F,I_M)
\right],
\qquad K=1,\ldots,T.
\label{eq:samcl-cumulative-objective}
\end{equation}
在顺序训练条件下，当前阶段通常只能完整访问 $\mathcal{D}^{(K)}$。若直接在新任务上继续优化，参数可能偏向当前分布，并降低旧任务上的解剖重叠或地标对应精度。

为提供有限历史证据，SAMCL 采用容量受限的历史回放集合 $\mathcal{M}$ 保存少量历史图像对，满足
$|\mathcal{M}|\leq b$。这一访问条件介于联合训练与无回放学习之间：模型不能重新调用全部历史数据，但后续更新仍会读取回放集合中的历史原始图像对。有限回放用于提供旧任务梯度，元学习用于减少当前梯度与历史梯度之间的干扰，锐度感知优化则关注参数邻域而非单个参数点的损失。最终模型既要保持已见任务的配准能力，也希望对尚未参与训练的任务形成更好的前向泛化，但这些目标只能在给定任务顺序、回放容量和评价数据上进行验证。

表\ref{tab:task-cl-problem-comparison}概括了本章两类研究的主要差别。该比较用于明确问题范围，而不是把 Dice、目标配准误差和形变质量压缩为同一数值。

\begin{table}[htbp]
\centering
\caption{本章弱监督持续分割与持续配准问题的主要差异}
\label{tab:task-cl-problem-comparison}
\thesistablefont
\renewcommand{\arraystretch}{1.18}
\setlength{\tabcolsep}{3pt}
\begin{tabularx}{\textwidth}{@{}p{1.9cm}XXXX@{}}
\toprule
研究对象 & 输出空间 & 当前阶段监督 & 历史访问 & 主要评价问题 \\
\midrule
ScribbleCL &
像素或体素级离散语义标签 &
少量涂鸦位置具有直接类别标签，并以全局一致性和空间先验补充约束 &
容量受限的历史图像--涂鸦与参考特征回放；无回放方法作为比较参照 &
当前任务建模、历史知识保持、遗忘、前向泛化及最终全类别分割 \\
\addlinespace
SAMCL &
固定图像域到移动图像域的连续空间变换 &
按任务使用图像相似性、解剖标签或地标构造损失 &
容量受限的历史图像对经验回放 &
已见任务配准精度、后向能力、形变合理性与前向泛化 \\
\bottomrule
\end{tabularx}
\end{table}

\section{ScribbleCL：面向涂鸦监督的持续医学图像分割研究}
\label{sec:scribblecl}

ScribbleCL 研究稀疏涂鸦监督与任务顺序变化叠加时的医学图像分割。当前阶段只利用涂鸦标签进行优化，历史任务则通过容量受限的回放集合参与后续训练。当前任务的弱监督目标采用本文前期工作中的全局一致性损失和空间先验损失\cite{zhang2026zscribbleseg}，不引入其他损失项。在此基础上，本文将当前任务学习、历史特征保持和回放样本上的弱监督优化整合，并将这一组合方法记为 ZSDERpp。该方法使用同一总体目标处理 Domain-CL、Class-CL 和 Organ-CL；其中，Class-CL 在统一目标内部增加背景语义校正，以避免新增类别到达后旧类别被错误压入背景\cite{cermelli2020mib}。

\subsection{涂鸦监督下的全局一致性与空间先验}
\label{subsec:prior-scribble-foundation}

设医学图像为 $X:\Omega\rightarrow\mathbb{R}^{C}$，涂鸦标签为 $Y^{\mathrm{scr}}$，具有有效标签的位置集合为 $\Omega_L$。部分交叉熵只在涂鸦位置产生直接监督：
\begin{equation}
\mathcal{L}_{\mathrm{PCE}}
=
-\frac{1}{|\Omega_L|}
\sum_{i\in\Omega_L}
\sum_{k=0}^{K-1}
y_{ik}\log p_{\theta,k}(x_i),
\label{eq:scribblecl-pce}
\end{equation}
其中，$K$ 为类别数，$y_{ik}$ 为像素 $x_i$ 的涂鸦标签，$p_{\theta,k}(x_i)$ 为模型预测概率。该损失不把未标注像素直接视为背景，但仅依赖少量涂鸦位置时，仍难以约束完整区域和轮廓。

\paragraph{全局一致性损失。}
全局一致性要求图像变换与预测变换具有相容结果\cite{zhang2026zscribbleseg}。设 $(X_1,X_2)$ 为同一批次中的两个输入，$\mathcal{A}(\cdot,\cdot)$ 表示对两幅图像或两组预测执行相同参数的混合与遮挡变换，定义
\begin{equation}
\begin{aligned}
u_{12}&=\mathcal{A}\!\left(f_{\theta}(X_1),f_{\theta}(X_2)\right),
&v_{12}&=f_{\theta}\!\left(\mathcal{A}(X_1,X_2)\right),\\
u_{21}&=\mathcal{A}\!\left(f_{\theta}(X_2),f_{\theta}(X_1)\right),
&v_{21}&=f_{\theta}\!\left(\mathcal{A}(X_2,X_1)\right).
\end{aligned}
\label{eq:scribblecl-global-pairs}
\end{equation}
令
$\ell_{\cos}(u,v)=-\langle u,v\rangle/(\|u\|_2\|v\|_2)$，则
\begin{equation}
\mathcal{L}_{\mathrm{global}}
=
\frac{1}{2}
\left[
\ell_{\cos}(u_{12},v_{12})
+
\ell_{\cos}(u_{21},v_{21})
\right].
\label{eq:scribblecl-global-loss}
\end{equation}
该项约束“先预测后变换”和“先变换后预测”的输出关系，使未标注区域能够通过增强前后的整体响应获得附加约束，而不被写成确定伪标签\cite{zhang2026zscribbleseg}。

\paragraph{空间先验损失。}
空间先验根据局部位置、图像外观和当前预测，从未标注区域中筛选相对可信的类别负集合。令 $\overline{\Omega}_k$ 表示对前景类别 $k$ 得到的负候选位置，将除 $k$ 外的类别合并为 $\overline c_k$，则空间损失写为
\begin{equation}
\mathcal{L}_{\mathrm{spatial}}
=
-
\sum_{k=1}^{K-1}
\sum_{x_i\in\overline{\Omega}_k}
\log p_{\theta}(\overline c_k\mid x_i).
\label{eq:scribblecl-spatial-loss}
\end{equation}
由此，该项利用相对可信的负证据限制类别区域不合理扩张。候选正区域不被直接当作无噪声标签，因此空间约束主要用于排除不符合局部结构的类别归属\cite{zhang2026zscribbleseg}。

本章当前任务只组合上述两项弱监督约束：
\begin{equation}
\mathcal{L}_{\mathrm{cur}}^{(t)}
=
\mathcal{L}_{\mathrm{PCE}}^{(t)}
+
\lambda_{g}\mathcal{L}_{\mathrm{global}}^{(t)}
+
\lambda_{s}\mathcal{L}_{\mathrm{spatial}}^{(t)}.
\label{eq:scribblecl-current-loss}
\end{equation}
其中，$\lambda_g$ 和 $\lambda_s$ 分别控制全局一致性与空间先验的强度。本章采用的弱监督目标由部分交叉熵、全局一致性和空间先验构成\cite{zhang2026zscribbleseg}。

\subsection{面向三类持续分割设定的 ZSDERpp}
\label{subsec:scribblecl-methods}

在历史信息组织上，设第 $t$ 个阶段的涂鸦训练集为 $\mathcal{D}^{\mathrm{scr}}_t$，历史回放集合为 $\mathcal{M}_{t-1}$，其容量固定为 64 幅图像，即 $|\mathcal{M}_{t-1}|\leq 64$。回放集合采用水库采样在线更新，使顺序到达的历史样本具有被保留的机会。每个回放条目包含历史图像、对应涂鸦标签、样本进入回放集合时保存的参考特征以及任务标识，可写为
$(X_m,Y_m^{\mathrm{scr}},\widetilde{\mathbf h}_m,\tau_m)$。任务标识用于在 Organ-CL 中选择对应输出头，在 Class-CL 中解释样本所处阶段的标签空间。

令 $\mathbf h_{\theta}(X)$ 表示实现中选定的中间特征，历史特征保持损失为
\begin{equation}
\mathcal{L}_{\mathrm{feature}}
=
\frac{1}{|\mathcal{B}_t|}
\sum_{m\in\mathcal{B}_t}
\left\|
\mathbf h_{\theta}(X_m)-\widetilde{\mathbf h}_m
\right\|_2^2,
\label{eq:scribblecl-feature-loss}
\end{equation}
其中，$\mathcal{B}_t\subseteq\mathcal{M}_{t-1}$ 为当前更新采样的历史小批次，回放批大小固定为 8，即 $|\mathcal{B}_t|=8$。该项限制历史样本的表征漂移；回放样本的涂鸦标签仍用于重新计算部分交叉熵、全局一致性和空间先验。历史回放同时利用参考特征和稀疏涂鸦标签，以约束旧任务表征并重新优化回放样本。

按照本章实验实现，回放损失固定为
\begin{equation}
\mathcal{L}_{\mathrm{replay}}^{(t)}
=
0.5\,\mathcal{L}_{\mathrm{feature}}^{(t)}
+
0.5\,\mathcal{L}_{\mathrm{PCE}}^{\mathrm{replay},(t)}
+
1.0\,\mathcal{L}_{\mathrm{global}}^{\mathrm{replay},(t)}
+
0.1\,\mathcal{L}_{\mathrm{spatial}}^{\mathrm{replay},(t)}.
\label{eq:scribblecl-replay-loss}
\end{equation}
由此，ZSDERpp 的统一训练目标为
\begin{equation}
\mathcal{L}_{\mathrm{ZSDERpp}}^{(t)}
=
\mathcal{L}_{\mathrm{cur}}^{(t)}
+
\mathbb{I}[t>1]\mathcal{L}_{\mathrm{replay}}^{(t)},
\label{eq:scribblecl-zsderpp-total}
\end{equation}
其中，第一阶段回放集合为空，历史损失取零。式\eqref{eq:scribblecl-replay-loss}把历史表征保持与回放样本上的稀疏监督重新优化置于同一目标中：特征项保持历史响应，部分交叉熵保留已标注像素的类别信息，全局一致性与空间先验则继续约束历史样本的未标注区域。

\paragraph{三类持续设定的统一适配。}
在 Domain-CL 中，标签语义和共享输出头保持不变，回放样本直接通过同一输出空间计算式\eqref{eq:scribblecl-replay-loss}。Organ-CL 为各任务保留独立输出头；历史条目的任务标识决定回放时使用的输出头，而共享编码器继续受到特征保持和弱监督回放约束。因此，两类设定不需要改变 ZSDERpp 的总体形式。

相比之下，Class-CL 逐阶段扩展统一输出空间，当前标签中的背景还可能包含旧类别。为避免把未标注旧类别强制压入背景，当前阶段背景标签的监督概率定义为\cite{cermelli2020mib}
\begin{equation}
\overline{p}_{t}(y\mid\mathbf{u})
=
\begin{cases}
p_t(y\mid\mathbf{u}), & y\in\mathcal{C}_t,\\[1mm]
p_t(b\mid\mathbf{u})+
\displaystyle\sum_{c\in\mathcal{C}_{1:t-1}}p_t(c\mid\mathbf{u}), & y=b,
\end{cases}
\label{eq:scribblecl-background-aware-probability}
\end{equation}
并以 $\overline p_t$ 计算背景感知的部分交叉熵
\begin{equation}
\mathcal{L}_{\mathrm{uPCE}}^{(t)}
=
-\frac{1}{|\Omega_L^{(t)}|}
\sum_{\mathbf{u}\in\Omega_L^{(t)}}
\log \overline p_t\!\left(y_{\mathbf{u}}^{(t)}\mid\mathbf{u}\right).
\label{eq:scribblecl-background-aware-pce}
\end{equation}
旧模型提供历史类别的函数约束。对旧类别保持逐类匹配，并将旧模型背景概率与当前模型的背景及新增类别概率之和对应：
\begin{equation}
\widetilde p_t(c\mid\mathbf{u})
=
\begin{cases}
p_t(c\mid\mathbf{u}), & c\in\mathcal{C}_{1:t-1},\\[1mm]
p_t(b\mid\mathbf{u})+
\displaystyle\sum_{k\in\mathcal{C}_t}p_t(k\mid\mathbf{u}), & c=b.
\end{cases}
\label{eq:scribblecl-background-aware-distillation-probability}
\end{equation}
相应的背景语义保持项为
\begin{equation}
\mathcal{L}_{\mathrm{bg}}^{(t)}
=
-\frac{1}{|\Omega|}
\sum_{\mathbf{u}\in\Omega}
\sum_{c\in\mathcal{C}_{1:t-1}\cup\{b\}}
p_{t-1}(c\mid\mathbf{u})
\log \widetilde p_t(c\mid\mathbf{u}).
\label{eq:scribblecl-background-aware-distillation}
\end{equation}
在 Class-CL 中，式\eqref{eq:scribblecl-current-loss}的部分交叉熵由 $\mathcal{L}_{\mathrm{uPCE}}^{(t)}$ 替代，并加入 $\lambda_{\mathrm{bg}}\mathcal{L}_{\mathrm{bg}}^{(t)}$；回放监督在统一标签空间中采用相同的背景语义解释。该调整属于 ZSDERpp 在类别增量输出空间中的内部组成，不另行定义第二种方法。

\begin{table}[htbp]
\centering
\caption{ScribbleCL 比较方法及其历史信息访问条件}
\label{tab:scribblecl-methods}
\thesistablefont
\renewcommand{\arraystretch}{1.15}
\setlength{\tabcolsep}{3pt}
\begin{tabularx}{\textwidth}{@{}p{2.8cm}Xp{4.1cm}X@{}}
\toprule
方法 & 当前阶段监督 & 历史知识处理 & 本章用途 \\
\midrule
Dense-Sequential & 密集标签 & 无 & 密集监督顺序训练参照 \\
PCE-Sequential & 涂鸦与部分交叉熵 & 无 & 基础弱监督参照 \\
ZS-Sequential & 涂鸦、全局一致性与空间先验 & 无 & 增强弱监督顺序参照 \\
ZS-EWC & 与 ZS-Sequential 相同 & 参数重要性正则 & 无原始样本回放比较 \\
ZS-GPM & 与 ZS-Sequential 相同 & 梯度子空间投影 & 无原始样本回放比较 \\
ZSDERpp & 与 ZS-Sequential 相同 & 64 幅图像；水库采样；回放批大小为 8；保存涂鸦与参考特征 & Domain-CL、Class-CL、Organ-CL \\
\bottomrule
\end{tabularx}
\end{table}

\subsection{实验设定}
\label{subsec:scribblecl-experimental-settings}

\subsubsection{数据集、任务序列与数据划分}
\label{subsubsec:scribblecl-datasets}

ScribbleCL 直接复用第\ref{chap:benchmark}章表~\ref{tab:benchmark-datasets}中的三个数据场景、任务顺序和病例级划分\cite{wang2026benchmark}。每个任务仍按 60\% 训练集、15\% 验证集和 25\% 测试集组织，不重新抽样，也不调整任务到达次序。唯一的数据标注变化发生在训练集：第三章使用阶段可访问的密集标签，本章从同一批密集标签生成固定涂鸦；验证集和测试集继续使用相同的密集标签。由此，Dense-Sequential、各弱监督方法以及第三章对应基准共享同一组训练、验证和测试病例。

\begin{table}[htbp]
\centering
\caption{ScribbleCL 对第三章三类基准数据与任务序列的沿用关系}
\label{tab:scribblecl-datasets}
\thesistablefont
\renewcommand{\arraystretch}{1.16}
\setlength{\tabcolsep}{3pt}
\begin{tabularx}{\textwidth}{@{}p{1.85cm}Xp{2.5cm}p{2.7cm}@{}}
\toprule
场景 & 数据集与任务顺序 & 输出组织 & 训练/验证/测试 \\
\midrule
Domain-CL &
NCI-ISBI 2013 的 A、B 中心，I2CVB 的 C 中心，以及 PROMISE12 的 D、E、F 中心，按 A$\rightarrow$B$\rightarrow$C$\rightarrow$D$\rightarrow$E$\rightarrow$F 到达 &
共享二通道输出头 &
60\%/15\%/25\% \\
\addlinespace
Class-CL &
MMWHS 心脏 CT，依次学习 LV、LA、MYO，第二阶段加入 RV、RA，第三阶段加入 AA、PA &
共享八通道输出头 &
60\%/15\%/25\% \\
\addlinespace
Organ-CL &
LAScarQS 左心房 LGE MRI $\rightarrow$ PROMISE12 中心 D 前列腺 MRI $\rightarrow$ LiTS 肝脏 CT $\rightarrow$ FeTS 2021 脑肿瘤 FLAIR MRI &
每个任务独立二通道输出头 &
60\%/15\%/25\% \\
\bottomrule
\end{tabularx}
\end{table}

Domain-CL 的 C 中心执行裁剪和空间对齐，所有轴位切片调整为 $256\times256$，并使用非零体素的均值和标准差进行归一化。Class-CL 将 MMWHS 图像的平面分辨率重采样为 $0.78\times0.78$ mm，平均层厚为 1.60 mm，随后调整为 $256\times256$。Organ-CL 对每幅图像进行零均值、单位方差归一化，并将轴位切片调整为 $256\times256$。这些处理、输出头设置和任务顺序均与第三章一致。表~\ref{tab:scribblecl-datasets}用于说明继承关系，不表示本章重新建立了独立的数据划分。

\subsubsection{参考标注、金标准与数据使用边界}
\label{subsubsec:scribblecl-reference-standard}

本节的分割金标准是各公开数据集随影像提供的密集参考标签，而不是模型生成的伪标签，也不是从涂鸦反推的完整掩膜。Domain-CL 沿用 NCI-ISBI 2013、I2CVB 和 PROMISE12 的前列腺参考分割\cite{dataset_prostate2,dataset_prostate,dataset_prostate3}；Class-CL 沿用 MMWHS 的七类心脏结构标签，分别为左心室、左心房、心肌、右心室、右心房、升主动脉和肺动脉\cite{care25ASA1,discussion_inter_seg3,care25ASA2}；Organ-CL 则沿用各任务对应的左心房、前列腺、肝脏和脑肿瘤参考区域。各来源采用各自的标注与质量控制流程，本文没有重新组织一套跨数据集、多标注者共识标注。这里的“金标准”因此指实验中固定的评价参照，不意味着所有器官边界均不存在标注分歧。

参考标签在不同数据子集中承担不同作用。训练集的密集标签用于离线构建模拟涂鸦，并供 Dense-Sequential 密集监督参照使用；弱监督方法训练时读取固定涂鸦及其忽略区域，不把完整训练掩膜加入损失。验证集保留密集标签，用于模型选择；测试集保留密集标签，只用于独立评价。因此，本节研究的是训练阶段监督受限，而不是验证与测试阶段也只拥有涂鸦。相同病例的切片随病例级划分归属同一子集，避免把切片数量误当成独立病例数量。

标签语义也随场景明确限定。Domain-CL 共享前列腺与背景两类语义。Class-CL 在同一八通道输出中保留背景和七类心脏结构，但训练阶段只开放当前协议允许的类别，其余区域按阶段背景规则处理；这种训练标签重映射不能改变测试参考标签的原始解剖含义。Organ-CL 分别使用任务对应的二通道输出头，依靠任务标识确定当前前景。各方法在同一测试标签上评价，才能把模型变化与参考标签变化区分开来。

涂鸦是由训练参考掩膜模拟生成的，因此其类别正确性继承自参考掩膜，而其位置、长度和形态由生成过程决定。该设计控制了不同方法之间的监督差异，却没有测量医生实际绘制涂鸦的时间，也没有覆盖漏标、错标和不同标注者习惯。现有结果不能据此换算为临床标注工时节约或真实标注者之间的一致性；这些问题需要额外的标注研究。

\subsubsection{涂鸦数据集构建}
\label{subsubsec:scribblecl-scribble-construction}

训练集涂鸦采用 WSL4MIS 开源实现中的 \texttt{scribbles\_generator.py} 构建\cite{luo2022scribble}\footnote{\url{https://github.com/HiLab-git/WSL4MIS/blob/main/code/scribbles_generator.py}}。对每个类别，先从阶段可访问的密集标签中得到二值掩膜，再逐轴位切片生成骨架。面积较大的区域先进行随机腐蚀，随后使用 Lee 骨架化获得中心线。前景骨架经过连通区域筛选和随机分支裁剪，并进行一次形态学膨胀。对面积超过 1000 像素的目标，脚本在原类别区域内施加 $[-6,6]$ 像素的随机平移和不超过 $15^\circ$ 的随机旋转，以增加涂鸦形态变化。

生成结果保留前景类别和背景涂鸦，其他位置写为 255，并在部分交叉熵中作为忽略位置。若训练代码内部使用其他 \texttt{ignore\_index}，只在数据读取阶段执行确定性映射，不改变磁盘中涂鸦的含义。所有比较方法共享同一份固定涂鸦文件，避免重新随机生成标注造成额外差异。

图\ref{fig:scribblecl-annotation-sequences}给出了三类持续分割场景中由密集训练标签生成的代表性涂鸦及任务到达顺序。Domain-CL 保持前列腺分割语义不变并依次改变数据中心，Organ-CL 依次改变器官与成像模态，Class-CL 则分阶段扩展心脏结构类别。

\begin{figure}[htbp]
\centering
\includegraphics[width=\textwidth]{ch04/scribble_annotation_fig.pdf}
\caption{ScribbleCL 三类持续分割场景的涂鸦标注与任务序列。每一行对应一种持续分割设定，箭头表示任务到达顺序；蓝色曲线表示背景涂鸦，其他彩色曲线表示当前阶段可访问的前景类别涂鸦，未被涂鸦覆盖的像素在弱监督训练中设为忽略位置。}
\label{fig:scribblecl-annotation-sequences}
\end{figure}

针对 Class-CL，先按照当前阶段的标签空间完成类别重映射，再调用涂鸦生成代码。当前不可访问的旧类别和未来类别按该阶段的背景语义处理。完整密集标签只用于生成模拟涂鸦、验证、测试和 Dense-Sequential；其他方法的当前任务损失与历史回放均只使用涂鸦标签，回放条目还保存由当时模型计算的参考特征。

\subsubsection{训练与评价设置}
\label{subsubsec:scribblecl-training-evaluation}

网络与基础训练预算沿用第三章：所有实验采用 ResUNet32，每个任务训练 150 个 epoch，优化器为 SGD，初始学习率为 0.008，在第 80 个 epoch 后乘以 0.5，批大小为 8。Dense-Sequential、PCE-Sequential、ZS-Sequential、ZS-EWC、ZS-GPM 和 ZSDERpp 使用相同的数据划分、任务顺序、网络与训练预算。当前任务损失中的权重设置为 $\lambda_g=0.05$ 和 $\lambda_s=1$；前 15 个 epoch 不启用空间先验，自第 16 个 epoch 起加入 $\mathcal{L}_{\mathrm{spatial}}$。上述设置在所有使用全局一致性与空间先验的比较方法中保持一致\cite{zhang2026zscribbleseg}。

在持续学习参数方面，ZS-EWC 和 ZS-GPM 的持续学习参数沿用第三章相应基线设置\cite{kirkpatrick2017ewc,saha2021gpm}。ZSDERpp 在每个后续阶段从容量为 64 幅图像的历史回放集合中采样图像--涂鸦对及参考特征，并采用水库采样更新该集合。回放批大小为 8，随机种子固定为 42。回放损失严格采用式\eqref{eq:scribblecl-replay-loss}中的 $0.5/0.5/1.0/0.1$ 系数。Class-CL 额外启用式\eqref{eq:scribblecl-background-aware-probability}--\eqref{eq:scribblecl-background-aware-distillation}所示的背景语义校正；Domain-CL 和 Organ-CL 不使用该调整。

评价直接采用第三章定义的阶段--任务性能矩阵及指标公式。A-Dice、BWTR 和 RMA 分别按式\eqref{eq:benchmark-adice}、式\eqref{eq:benchmark-bwtr}和式\eqref{eq:benchmark-rma}计算；Domain-CL 进一步使用式\eqref{eq:benchmark-efwt}中的 E-FWT，Class-CL 使用式\eqref{eq:benchmark-wcd}中的 WCD。Organ-CL 报告 A-Dice、BWTR 和 RMA。所有指标均在与第三章相同的密集测试标签上计算，指标方向和解释也保持一致。这样，第三章与第四章的结果可以在同一评价维度下讨论，并分析监督形式与历史信息使用方式对各项能力的影响。鉴于现有定量证据仅覆盖 Domain-CL，结果分析限定于该场景；Class-CL 与 Organ-CL 保留方法和实验设置，不作性能判断。

\subsection{实验结果}
\label{subsec:scribblecl-results}

表\ref{tab:scribblecl-domain-results}给出了六中心前列腺 MRI 序列上的 Domain-CL 结果。A-Dice、RMA 和 E-FWT 的数值越大越好；BWTR 数值越大，表示后续阶段训练后旧域性能保持越充分。

\begin{table}[htbp]
\centering
\caption{Domain-CL 场景下 ScribbleCL 方法的结果。红色粗体表示各指标的最高点估计；不同监督与回放条件并不等价。}
\label{tab:scribblecl-domain-results}
\thesistablefont
\renewcommand{\arraystretch}{1.14}
\setlength{\tabcolsep}{3pt}
\begin{tabular}{lcccc}
\toprule
方法 & A-Dice & BWTR & RMA & E-FWT \\
\midrule
Dense-Sequential & $0.676$ & $-0.237$ & $\best{1.086}$ & $\best{0.260}$ \\
PCE-Sequential & $0.248$ & $-0.545$ & $0.776$ & $0.133$ \\
ZS-Sequential & $0.510$ & $-0.324$ & $0.732$ & $0.234$ \\
ZS-EWC & $0.543$ & $-0.298$ & $0.761$ & $0.216$ \\
ZS-GPM & $0.612$ & $-0.206$ & $0.808$ & $0.241$ \\
ZSDERpp & $\best{0.701}$ & $\best{-0.152}$ & $0.899$ & $0.225$ \\
\bottomrule
\end{tabular}
\end{table}

\subsection{实验分析}
\label{subsec:scribblecl-experimental-analysis}

Dense-Sequential 的 A-Dice、BWTR、RMA 和 E-FWT 分别为 $0.676$、$-0.237$、$1.086$ 和 $0.260$。其 RMA 与 E-FWT 为表中最高值，说明密集监督直接顺序训练在当前域相对建模和前向泛化方面仍具有优势；BWTR 为负则表明，密集标签本身不能消除六中心顺序更新中的旧域性能下降。

在稀疏监督参照中，PCE-Sequential 只在涂鸦位置计算部分交叉熵。与其相比，ZS-Sequential 的 A-Dice 提高 $0.262$，BWTR 提高 $0.221$，E-FWT 提高 $0.101$，而 RMA 降低 $0.044$。在相同涂鸦、骨干和训练预算下，全局一致性与空间先验改善了最终多域性能、旧域保持和前向泛化表现，但没有同步提高当前任务相对于独立训练参照的建模能力。这一差异说明，弱监督持续分割仍需联合观察主对角线、历史任务和未来任务条目。

进一步加入持续学习约束后，两种无原始样本回放的持续学习比较方法表现出不同取舍。ZS-EWC 相较 ZS-Sequential 将 A-Dice、BWTR 和 RMA 分别提高 $0.033$、$0.026$ 和 $0.029$，但 E-FWT 降低 $0.018$。ZS-GPM 的对应变化为 $+0.102$、$+0.118$、$+0.076$ 和 $+0.007$，在四项指标上均高于 ZS-Sequential，并在无回放弱监督方法中取得最高 A-Dice、BWTR、RMA 和 E-FWT。在当前 Domain-CL 序列中，梯度方向保护呈现出比单纯参数重要性约束更平衡的稳定性--可塑性结果。

在有限回放条件下，ZSDERpp 的 A-Dice 和 BWTR 分别为 $0.701$ 和 $-0.152$，均为表中最高值。与 ZS-Sequential 相比，两项指标分别提高 $0.191$ 和 $0.172$，RMA 提高 $0.167$，E-FWT 降低 $0.009$；与 ZS-GPM 相比，A-Dice、BWTR 和 RMA 分别提高 $0.089$、$0.054$ 和 $0.091$，E-FWT 降低 $0.016$。有限历史回放、特征保持以及回放样本上的弱监督重新优化共同提高了最终多域性能和旧域保持，但更强的历史约束没有带来最高的前向泛化。

相较 Dense-Sequential，ZSDERpp 的 A-Dice 高 $0.025$，BWTR 高 $0.085$，但 RMA 和 E-FWT 分别低 $0.187$ 和 $0.035$。在当前回放容量下，有限历史回放提高了 ZSDERpp 的最终平均性能和知识保持，但当前任务相对建模与前向泛化仍低于密集监督参照。

\subsection{讨论}
\label{subsec:scribblecl-discussion}

本章采用本文前期工作中的全局一致性和空间先验\cite{zhang2026zscribbleseg}。前者利用增强前后的输出关系补充未标注区域约束，后者用空间负证据限制类别区域扩张。PCE-Sequential 与 ZS-Sequential 的比较表明，这两项约束在 Domain-CL 中提高了 A-Dice、BWTR 和 E-FWT；RMA 的下降也说明，增强弱监督目标并不必然提高每个新域的独立建模能力。

ZSDERpp 保存容量受限的历史图像--涂鸦对和参考特征，并通过式\eqref{eq:scribblecl-replay-loss}重新约束历史表征、已标注像素和未标注区域。有限历史回放使模型能够在后续阶段重新优化旧域样本，从而改善 A-Dice 与 BWTR。

Class-CL 与 Organ-CL 使用同一 ZSDERpp 总体目标，但输出组织不同。Class-CL 需要在当前监督和旧模型约束中重新解释背景概率，Organ-CL 则依靠任务标识路由至对应输出头。三类设定共享总体训练框架；值得说明的是，现有定量证据不足以判断 Class-CL 与 Organ-CL 的实际性能。

本节实验采用 A 到 F 的六中心前列腺 MRI 顺序、固定数据划分、由密集训练标签模拟得到的涂鸦、ResUNet32 骨干、统一的 150-epoch 训练预算和当前回放配置。模拟涂鸦不能覆盖真实标注者在位置、长度和类别选择上的差异。

\section{SAMCL：锐度感知元持续医学图像配准方法}
\label{sec:samcl}

\subsection{研究动机：任务扩展中的配准更新与知识保持}
\label{subsec:samcl-motivation}

医学图像配准通过估计固定图像与移动图像之间的空间变换，使两幅图像中的对应解剖结构在同一坐标系下对齐，是疾病进展监测、术前规划、手术导航和多模态影像联合分析中的基础环节\cite{maintz1998registration,gupta2021registration,maes1997mutualinformation,sotiras2013deformable}。深度学习方法能够直接从图像对学习配准变换，在推理效率和任务内配准精度方面取得了较好结果\cite{haskins2020deepreg,balakrishnan2019voxelmorph}。然而，现有网络通常围绕特定解剖区域、成像模态和配准关系构建，例如脑部 MR--MR、肺部 CT--CT 或腹部 CT--CT 配准\cite{kang2022dual,che2023amnet,hering2021lungreg}。一旦测试数据的器官、外观、模态关系或采集条件偏离训练分布，任务特定模型的性能容易下降。

然而，真实临床数据还具有持续到达的特点。新设备投入使用、检查方案变化以及新的器官或模态进入分析流程，都会使配准任务随时间扩展。若为每个任务分别保存一套网络，模型数量和维护成本会随任务增加；若等待全部数据集中后再进行多任务训练，又要求不同阶段的数据能够被长期保存并集中访问，这在医疗数据治理和隐私要求下往往难以实现\cite{lee2020clinicalcl,vokinger2021medicaldevices}。直接使用当前任务继续更新同一网络虽然满足顺序训练条件，却会因分布变化破坏早期任务的参数解，产生灾难性遗忘\cite{thrun1995lifelong,mccloskey1989catastrophic,delange2022continual}。

基于上述观察，本文将跨器官与跨模态配准组织为持续学习问题：模型依次接收来自不同器官、模态和配准关系的数据集，在不能重新集中访问全部历史数据的条件下持续更新；训练结束后，单一模型应尽可能保持所有已见任务的配准能力，并对后续未见任务具有一定前向泛化。这里的“通用”并不表示对任意临床数据都具有无条件保证，而是指在给定任务序列和访问约束下，以一个持续更新的模型覆盖多项临床相关配准任务\cite{wang2024samcl}。

图\ref{fig:samcl-framework}展示了这一设定。训练阶段依次到达脑部 MR--MR、腹部 CT--CT、肺部 CT--CT 和腹部 MR--CT 图像对，任务之间同时存在解剖区域、图像外观、模态关系和个体内/个体间配准方式的变化。推理阶段不再为各任务维护独立网络，而是使用顺序训练后得到的单一模型处理已见任务。SAMCL 以有限经验回放保存少量历史图像对，以一阶元持续学习协调当前和历史样本的更新方向，并进一步通过锐度感知优化改善参数邻域内的损失，从而同时面向历史保持和前向泛化。

\IfFileExists{figures/ch04/samcl_framework.pdf}{%
\begin{figure}[htbp]
\centering
\includegraphics[width=\textwidth]{ch04/samcl_framework.pdf}
\caption[SAMCL 的持续医学图像配准设定]{SAMCL 的持续医学图像配准设定。训练阶段依次接收脑部 MR--MR、腹部 CT--CT、肺部 CT--CT 和腹部 MR--CT 图像对；不同任务在解剖区域、成像模态和配准关系上存在明显分布差异。完成顺序训练后，以一个模型处理已见任务，用于已见任务的配准推理；图示不代表真实临床系统部署。}
\label{fig:samcl-framework}
\end{figure}
}{}

\subsection{配准网络与持续学习目标}
\label{subsec:samcl-registration-objective}

设固定图像为 $f:\Omega_F\rightarrow\mathbb{R}$，移动图像为 $m:\Omega_M\rightarrow\mathbb{R}$。医学图像配准需要估计从固定图像域到移动图像域的空间变换
$\phi:\Omega_F\rightarrow\Omega_M$，使变形后的移动图像
$m\circ\phi$ 与固定图像中的对应结构对齐。学习式配准网络以图像对为输入，并用参数 $\boldsymbol{\theta}$ 表示该变换：
\begin{equation}
\phi_{\boldsymbol{\theta}}
=
g(f,m;\boldsymbol{\theta}).
\label{eq:samcl-registration-network}
\end{equation}
对单一数据集 $\mathcal{D}=\{(f_i,m_i)\}_{i=1}^{N}$，模型通过最小化经验风险学习参数：
\begin{equation}
\widehat{\boldsymbol{\theta}}
=
\operatorname*{arg\,min}_{\boldsymbol{\theta}}
\mathbb{E}_{(f,m)\sim\mathcal{D}}
\left[
\mathcal{L}(\phi_{\boldsymbol{\theta}};f,m)
\right],
\label{eq:samcl-single-task-objective}
\end{equation}
其中，配准损失可统一写为
\begin{equation}
\mathcal{L}(\phi_{\boldsymbol{\theta}};f,m)
=
D\!\left(f,m\circ\phi_{\boldsymbol{\theta}}\right)
+
R\!\left(\phi_{\boldsymbol{\theta}}\right).
\label{eq:samcl-registration-loss}
\end{equation}
$D(\cdot,\cdot)$ 衡量固定图像与变形后移动图像的外观或解剖不一致，$R(\cdot)$ 用于约束形变的平滑性。不同任务能够提供的标签或地标不同，因此第 $t$ 个任务使用任务特定损失 $\mathcal{L}^{(t)}$。

在持续配准中，数据集
$\{\mathcal{D}^{(t)}\}_{t=1}^{T}$ 按顺序到达。若完成前 $K$ 个任务后仍能访问所有已见数据，则理想目标为式\eqref{eq:samcl-cumulative-objective}中的累积风险最小化。实际训练阶段通常只能完整访问当前数据集 $\mathcal{D}^{(K)}$，因此需要在有限历史访问条件下近似这一目标。SAMCL 保存容量受限的历史回放集合 $\mathcal{M}$，后续阶段能够重新读取其中的固定--移动图像对，但不能访问完整历史训练集。该设定介于集中式多任务训练和无回放持续学习之间。

\subsection{有限经验回放下的元持续学习}
\label{subsec:samcl-meta-continual}

仅将当前任务损失与回放损失相加，不能保证两类更新方向相互兼容。当当前任务梯度与历史任务梯度方向相反时，一次有利于新任务的参数更新可能增大旧任务损失。元持续学习从过去知识保持和未来任务适应两个方向组织更新，通过一阶元学习鼓励不同样本或任务之间的迁移，并减小梯度干扰。SAMCL 以 Meta Experience Replay（MER）为基础，MER 结合经验回放与 Reptile 一阶元学习\cite{riemer2019mer,nichol2018firstorder,method_replay_er}。

具体地，令 $\mathcal{Q}_t$ 表示由当前任务数据与历史回放集合共同形成的采样分布，并从中得到样本序列
$z_i=(f_i,m_i)$。样本对应任务记为 $t_i$，内层学习率为 $\alpha$。MER 近似优化下式：
\begin{equation}
\begin{aligned}
\min_{\boldsymbol{\theta}}
\mathbb{E}_{z_1,\ldots,z_s\sim\mathcal{Q}_t}
\Bigg[
\sum_{i=1}^{s}
\Bigg(&
\mathcal{L}^{(t_i)}_i(\boldsymbol{\theta})
-
\frac{\alpha}{2}
\sum_{j=1}^{i-1}
\nabla_{\boldsymbol{\theta}}
\mathcal{L}^{(t_i)}_i(\boldsymbol{\theta})^{\top}
\nabla_{\boldsymbol{\theta}}
\mathcal{L}^{(t_j)}_j(\boldsymbol{\theta})
\Bigg)
\Bigg].
\end{aligned}
\label{eq:samcl-meta-objective}
\end{equation}
第一项降低各样本自身的配准损失；第二项显式体现梯度内积。若两个梯度的内积为正，沿其中一个样本的下降方向更新通常也有利于另一样本；若内积为负，则表明两次更新存在干扰。因此，一阶元更新不只追求当前批次损失下降，还通过训练轨迹近似增加跨任务梯度的一致性，使模型在有限记忆下兼顾新任务学习和历史任务保持。

为在固定容量下保留历史样本，SAMCL 使用水库采样更新历史回放集合。设到当前时刻已经观察到 $n$ 个图像对，回放容量为 $B_M$。当回放集合未满时，当前图像对直接写入；回放集合填满后，第 $n$ 个样本以 $B_M/n$ 的概率进入回放集合，并随机替换其中一个已有样本。由此，顺序数据流中的样本具有近似相同的保留机会，回放集合不需要预先知道任务长度。水库采样只控制历史样本的进入方式，并不能保证不同器官、模态或任务在有限容量下完全均衡，因此元更新的效果仍依赖回放集合对历史分布的代表性。

\subsection{锐度感知元持续学习}
\label{subsec:samcl-sharpness-update}

可变形配准损失通常高度非凸，可能包含多个局部极小值\cite{hill2001registration}。在有限样本和顺序分布变化下，仅在当前参数点取得较低损失，并不意味着其邻域内的损失同样稳定；如果参数附近的损失曲面较尖锐，轻微的任务变化、批次采样差异或后续更新都可能导致性能明显下降。为此，SAMCL 将锐度感知最小化引入元持续学习\cite{foret2021sam}。

对任务 $t$，锐度感知损失定义为
\begin{equation}
\mathcal{L}^{(t)}_{\mathrm{sam}}(\boldsymbol{\theta})
=
\max_{\|\boldsymbol{\epsilon}\|_2\leq\rho}
\mathcal{L}^{(t)}
\left(\phi_{\boldsymbol{\theta}+\boldsymbol{\epsilon}}\right),
\label{eq:samcl-sharpness-loss}
\end{equation}
其中，$\rho$ 为参数邻域半径，$\boldsymbol{\epsilon}$ 表示邻域内的参数扰动。该目标不是只降低 $\boldsymbol{\theta}$ 处的损失，而是降低半径为 $\rho$ 的邻域内近似最坏损失，从而偏向具有较平坦损失曲面的参数区域。

将式\eqref{eq:samcl-sharpness-loss}代入元持续学习目标，可得 SAMCL 的总体目标：
\begin{equation}
\begin{aligned}
\min_{\boldsymbol{\theta}}
\mathbb{E}_{z_1,\ldots,z_s\sim\mathcal{Q}_t}
\Bigg[
\sum_{i=1}^{s}
\Bigg(&
\mathcal{L}^{(t_i)}_{\mathrm{sam},i}(\boldsymbol{\theta})
-
\frac{\alpha}{2}
\sum_{j=1}^{i-1}
\nabla_{\boldsymbol{\theta}}
\mathcal{L}^{(t_i)}_{\mathrm{sam},i}(\boldsymbol{\theta})^{\top}
\nabla_{\boldsymbol{\theta}}
\mathcal{L}^{(t_j)}_{\mathrm{sam},j}(\boldsymbol{\theta})
\Bigg)
\Bigg].
\end{aligned}
\label{eq:samcl-full-objective}
\end{equation}
该式同时包含两个层面的约束：梯度内积项面向当前样本与历史样本之间的迁移和干扰，锐度感知损失则面向每个样本对应的局部损失几何。二者分别作用于跨样本更新关系和参数邻域稳定性。

内层最大化无法直接求解。对
$\mathcal{L}^{(t)}(\boldsymbol{\theta}+\boldsymbol{\epsilon})$
在 $\boldsymbol{\epsilon}=\boldsymbol{0}$ 处作一阶泰勒展开：
\begin{equation}
\mathcal{L}^{(t)}(\boldsymbol{\theta}+\boldsymbol{\epsilon})
\approx
\mathcal{L}^{(t)}(\boldsymbol{\theta})
+
\boldsymbol{\epsilon}^{\top}
\nabla_{\boldsymbol{\theta}}
\mathcal{L}^{(t)}(\boldsymbol{\theta}).
\label{eq:samcl-taylor}
\end{equation}
在 $\ell_2$ 邻域约束下，上式的最大化可写为对偶范数问题，其近似最优扰动为
\begin{equation}
\widehat{\boldsymbol{\epsilon}}(\boldsymbol{\theta})
=
\rho
\frac{
\nabla_{\boldsymbol{\theta}}
\mathcal{L}^{(t)}(\boldsymbol{\theta})
}{
\left\|
\nabla_{\boldsymbol{\theta}}
\mathcal{L}^{(t)}(\boldsymbol{\theta})
\right\|_2
}.
\label{eq:samcl-perturbation}
\end{equation}
实际更新先在当前参数处计算梯度并构造
$\widehat{\boldsymbol{\epsilon}}$，再在扰动参数
$\boldsymbol{\theta}+\widehat{\boldsymbol{\epsilon}}$
处计算用于下降的梯度，即
$\nabla_{\boldsymbol{\theta}}\mathcal{L}^{(t)}_{\mathrm{sam}}$
由扰动点的梯度进行近似。SAMCL 对当前样本和历史回放样本均采用这一锐度感知更新，再通过元更新协调两条训练轨迹。经验回放提供历史任务证据，元持续学习减小任务间梯度干扰，锐度感知优化则提高局部参数邻域的稳健性。

\subsection{训练流程与配准模型}
\label{subsec:samcl-algorithm}

算法\ref{alg:samcl}给出 SAMCL 的训练流程。第 $t$ 个数据集到达后，模型从上一任务参数
$\widehat{\boldsymbol{\theta}}^{(t-1)}$ 继续训练。对当前图像对，算法首先保存更新前参数 $\boldsymbol{\theta}_0$，随后根据式\eqref{eq:samcl-perturbation}执行当前任务的锐度感知内层更新；当前图像对通过水库采样写入历史回放集合。然后，从回放集合采样历史图像对，按照其所属任务损失再次执行锐度感知更新。最后使用元学习率 $\beta$，从更新前参数 $\boldsymbol{\theta}_0$ 向内层更新后的参数 $\boldsymbol{\theta}$ 插值，完成一阶元更新。任务结束时保存
$\widehat{\boldsymbol{\theta}}^{(t)}$，并作为下一任务的初始化。

\begin{algorithm}[htbp]
\caption{锐度感知元持续医学图像配准（SAMCL）}
\label{alg:samcl}
\KwData{顺序到达的训练集 $\mathcal{D}^{(t)}=\{(f_i,m_i)\}_{i=1}^{N_t}$，$t=1,\ldots,T$}
\KwIn{内层学习率 $\alpha$，元学习率 $\beta$，邻域半径 $\rho$，内层批大小 $s$}
\KwOut{完成全部任务后的参数 $\widehat{\boldsymbol{\theta}}^{(T)}$}
随机初始化 $\widehat{\boldsymbol{\theta}}^{(0)}$，令历史回放集合 $\mathcal{M}\leftarrow\varnothing$\;
\For{$t=1,\ldots,T$}{
  将 $\mathcal{D}^{(t)}$ 组织为当前任务批次 $\mathcal{B}_t=\{B_i\}_{i=1}^{b_t}$，令 $\boldsymbol{\theta}\leftarrow\widehat{\boldsymbol{\theta}}^{(t-1)}$\;
  \For{$i=1,\ldots,b_t$}{
    $\boldsymbol{\theta}_0\leftarrow\boldsymbol{\theta}$\;
    \For{$j=1,\ldots,s$}{
      $\boldsymbol{g}\leftarrow\nabla_{\boldsymbol{\theta}}\mathcal{L}^{(t)}(\boldsymbol{\theta};B_i[j])$，$\widehat{\boldsymbol{\epsilon}}\leftarrow\rho\boldsymbol{g}/\|\boldsymbol{g}\|_2$\;
      $\boldsymbol{\theta}\leftarrow\boldsymbol{\theta}-\alpha\nabla_{\boldsymbol{\theta}}\mathcal{L}^{(t)}(\boldsymbol{\theta}+\widehat{\boldsymbol{\epsilon}};B_i[j])$\;
    }
    使用水库采样将当前图像对写入 $\mathcal{M}$\;
    从 $\mathcal{M}$ 中采样历史批次 $B_i^{\mathcal{M}}$，其所属任务记为 $t_i$\;
    $\boldsymbol{g}_{\mathcal{M}}\leftarrow\nabla_{\boldsymbol{\theta}}\mathcal{L}^{(t_i)}(\boldsymbol{\theta};B_i^{\mathcal{M}})$，$\widehat{\boldsymbol{\epsilon}}_{\mathcal{M}}\leftarrow\rho\boldsymbol{g}_{\mathcal{M}}/\|\boldsymbol{g}_{\mathcal{M}}\|_2$\;
    $\boldsymbol{\theta}\leftarrow\boldsymbol{\theta}-\alpha\nabla_{\boldsymbol{\theta}}\mathcal{L}^{(t_i)}(\boldsymbol{\theta}+\widehat{\boldsymbol{\epsilon}}_{\mathcal{M}};B_i^{\mathcal{M}})$\;
    $\boldsymbol{\theta}\leftarrow\boldsymbol{\theta}_0+\beta(\boldsymbol{\theta}-\boldsymbol{\theta}_0)$\tcp*{一阶元更新}
  }
  $\widehat{\boldsymbol{\theta}}^{(t)}\leftarrow\boldsymbol{\theta}$\;
}
\Return{$\widehat{\boldsymbol{\theta}}^{(T)}$}\;
\end{algorithm}

配准骨干采用 Large Kernel U-Net（LK-UNet），首层包含 16 个卷积核\cite{jia2022lkunet}。网络输出平稳速度场，并通过微分同胚积分获得空间变换；固定图像和移动图像之间采用对称配准。形变正则项由膜能和弯曲能组成，用于约束局部位移变化和形变曲率\cite{ashburner2007diffeomorphic}。这些正则项控制形变平滑性，而 SAMCL 的元更新和锐度感知目标分别面向跨任务干扰与损失曲面，两者不能替代配准模型自身的形变约束。

\subsection{配准持续学习实验}
\label{subsec:samcl-experiments}

\subsubsection{实验材料与任务序列}
\label{subsubsec:samcl-materials}

实验采用 Learn2Reg 挑战中的四项三维医学图像配准任务，并按照 OASIS 脑部 MR--MR、腹部 CT--CT、NLST 肺部 CT--CT 和腹部 MR--CT 的顺序到达\cite{hering2021learn2reg}。OASIS 来源于开放获取脑部 MRI 数据集，用于脑部个体间配准\cite{marcus2007oasis}；腹部 CT--CT 任务用于腹部个体间配准\cite{xu2016abdomenreg}；NLST 肺部 CT--CT 与腹部 MR--CT 数据来自公开影像资源，用于个体内肺部配准和跨模态腹部配准\cite{clark2013tcia,xu2016abdomenreg}。任务序列同时改变了解剖区域、图像外观、成像模态以及个体内/个体间关系，从而构成较强的连续分布变化。

相应地，表\ref{tab:samcl-datasets}给出四项任务的数据组织。OASIS、腹部 CT--CT 和腹部 MR--CT 分别使用 35、13 和 9 个对应解剖标签评价配准后的结构重叠；NLST 使用对应地标评价目标配准误差。表中的训练、验证和测试数量与本文前期实验设置保持一致\cite{wang2024samcl}，计量单位为各任务用于训练或评价的三维图像数据。

\begin{table}[htbp]
\centering
\caption{SAMCL 四个顺序配准任务的数据组织、预处理与评价}
\label{tab:samcl-datasets}
\thesistablefont
\renewcommand{\arraystretch}{1.14}
\setlength{\tabcolsep}{3pt}
\begin{tabularx}{\textwidth}{@{}cXp{1.8cm}p{2.35cm}p{1.65cm}p{2.0cm}p{1.5cm}@{}}
\toprule
顺序 & 数据集与任务 & 配准关系 & 训练/验证/测试 & 重采样间距 & 标签或地标 & 评价指标 \\
\midrule
1 & OASIS，脑部 MR--MR & 个体间 & 290/40/84 & $2\times2\times2$ mm & 35 类解剖标签 & 平均 Dice \\
2 & 腹部 CT--CT & 个体间 & 21/3/6 & $3\times3\times3$ mm & 13 类解剖标签 & 平均 Dice \\
3 & NLST，肺部 CT--CT & 个体内 & 294/42/84 & $3\times3\times3$ mm & 对应地标 & TRE（mm） \\
4 & 腹部 MR--CT & 个体内、跨模态 & 80/9/16 & $3\times3\times3$ mm & 9 类解剖标签 & 平均 Dice \\
\bottomrule
\end{tabularx}
\end{table}

所有体数据均裁剪或填充为
$112\times96\times112$。OASIS 重采样到
$2\times2\times2$ mm，其余任务重采样到
$3\times3\times3$ mm。CT 图像强度裁剪到
$[-200,300]$，以覆盖腹部软组织和器官的主要强度范围；随后，CT 和 MR 图像均线性归一化到 $[0,1]$。为避免不同方法因数据增强策略不同而获得额外收益，训练阶段不使用数据增强。

\subsubsection{对比方法与历史访问条件}
\label{subsubsec:samcl-baselines}

实验从任务特定训练、集中式训练、直接顺序训练和持续学习四个层面设置对比方法。未配准输入用于给出固定图像与移动图像在变形前的基础差异；独立训练和集中式多任务训练分别作为任务特定参考与完整数据访问参考；直接顺序训练用于观察不采用抗遗忘机制时的性能变化。持续学习方法包括 EWC、SI、GPM 和 MER。

EWC 根据旧任务估计参数重要性，并通过 Fisher 信息加权的二次惩罚限制重要参数偏移\cite{kirkpatrick2017ewc}。SI 在训练过程中累计参数对历史损失下降的贡献，以在线重要性约束后续更新\cite{method_regu_si}。GPM 从历史任务的表示中构建重要子空间，将新任务梯度投影到其正交补，以减小对旧任务方向的干扰\cite{saha2021gpm}。上述三种方法不回放历史原始图像。MER 使用有限经验回放和一阶元学习，通过增加跨样本梯度对齐来兼顾迁移与遗忘\cite{riemer2019mer}。SAMCL 与 MER 使用相同类型的有限历史图像对访问条件，其差别在于 SAMCL 将锐度感知更新嵌入当前和回放轨迹。

为明确资源条件，表\ref{tab:samcl-baselines}汇总了各方法的训练组织和历史访问范围。独立训练为每项任务保存不同模型；集中式多任务训练在一个模型中同时访问全部任务数据。二者不满足持续学习中的顺序访问条件，仅用于提供任务特定上限和完整历史访问下的参考表现。

\begin{table}[htbp]
\centering
\caption{SAMCL 实验中的对比方法及其训练信息访问条件}
\label{tab:samcl-baselines}
\thesistablefont
\renewcommand{\arraystretch}{1.16}
\setlength{\tabcolsep}{3pt}
\begin{tabular}{p{1.8cm}p{4.6cm}p{3.2cm}p{2.8cm}}
\toprule
方法 & 训练组织 & 历史原始图像访问 & 模型组织 \\
\midrule
未配准 & 不估计空间变换，直接评价固定图像与移动图像 & 不涉及 & 不涉及 \\
独立训练 & 每个任务仅使用本任务数据单独训练 & 不跨任务访问 & 每任务一个模型 \\
集中式多任务 & 所有任务数据同时参与联合训练 & 完整访问全部任务数据 & 单一模型 \\
直接顺序训练 & 当前任务到达后直接继续优化 & 不访问历史图像 & 单一持续模型 \\
EWC & 以参数重要性正则限制重要参数变化 & 不访问历史图像 & 单一持续模型 \\
SI & 以在线累计的重要性约束参数变化 & 不访问历史图像 & 单一持续模型 \\
GPM & 通过历史重要子空间投影新任务梯度 & 不访问历史图像 & 单一持续模型 \\
MER & 有限经验回放与一阶元持续学习 & 容量受限的历史图像对 & 单一持续模型 \\
SAMCL & 有限经验回放、元持续学习与锐度感知更新 & 容量受限的历史图像对 & 单一持续模型 \\
\bottomrule
\end{tabular}
\end{table}

\subsubsection{评价指标与实现细节}
\label{subsubsec:samcl-implementation}

对于 OASIS、腹部 CT--CT 和腹部 MR--CT，实验在相应解剖标签上计算平均 Dice；对于 NLST，在对应地标上计算以毫米为单位的目标配准误差（target registration error, TRE）。令 $R_{i,j}$ 表示完成任务 $i$ 后在任务 $j$ 测试集上的评价，则任务 $j$ 在全部任务训练结束后的最终性能和后向变化定义为
\begin{equation}
\mathrm{AVG}_j=R_{T,j},
\qquad
\mathrm{BWT}_j=R_{T,j}-R_{j,j}.
\label{eq:samcl-avg-bwt}
\end{equation}
对 Dice，数值越高越好，负 BWT 表示后续任务导致性能下降；对 TRE，数值越低越好，正 BWT 表示误差增大。因此，Dice-BWT 越高或越接近零越好，TRE-BWT 越低或越接近零越好。最后一个任务之后没有新的训练阶段，其 BWT 按定义为零。由于 Dice 和 TRE 的方向及量纲不同，本节不将四项任务的原始指标直接求平均。

在实现方面，所有方法使用相同的 LK-UNet 骨干。实验在 PyTorch 中实现，并在 NVIDIA RTX 3090 GPU 上运行。每个任务训练 $10000$ 次迭代，采用 Adam 优化器\cite{kingma2015adam}，学习率为 $1\times10^{-4}$，批大小为 $4$。SAMCL 的历史回放容量设为 $200$，元学习率 $\beta$ 设为 $0.25$，锐度邻域半径 $\rho$ 设为 $0.05$；其中，$\rho=0.05$ 是实验中表现最好的设置。

四项任务的标注形式不同，因此训练损失随任务变化。OASIS、腹部 CT--CT 和腹部 MR--CT 的损失由窗口大小为 $3$ 的负局部归一化互相关（local normalized cross-correlation, LNCC）和软 Dice 损失组成\cite{avants2008syn,hu2018weakreg}；NLST 的损失由负 LNCC 和基于对应地标的距离损失组成。所有任务均加入膜能和弯曲能构成的形变正则项，两项正则权重均设为 $1$。配准在固定图像与移动图像之间对称执行，并使用由平稳速度场参数化的微分同胚变换。

\subsubsection{配准精度与知识保持分析}
\label{subsubsec:samcl-analysis}

表\ref{tab:samcl-main-results}给出四任务顺序训练结束后的最终性能和后向变化。直接顺序训练在前三个任务上出现明显知识损失：OASIS 和腹部 CT--CT 的 Dice-BWT 分别为 $-0.395$ 和 $-0.129$，NLST 的 TRE-BWT 为 $3.508$ mm。SAMCL 将三项后向变化改善为 $-0.037$、$-0.039$ 和 $0.306$ mm，说明有限回放、元持续更新与锐度感知训练的组合能够明显缓解直接顺序训练的遗忘。

\begin{table}[htbp]
\centering
\caption{SAMCL 四任务顺序实验的最终性能与后向变化。}
\label{tab:samcl-main-results}
\thesistablefont
\renewcommand{\arraystretch}{1.12}
\setlength{\tabcolsep}{3pt}
\begin{tabular}{@{}lcccc@{}}
\toprule
\multicolumn{5}{c}{（a）最终性能 $R_{T,j}$} \\
方法 & OASIS Dice & 腹部 CT Dice & NLST TRE & 腹部 MR--CT Dice \\
\midrule
未配准 & $\underline{0.598}^{*}$ & $\underline{0.311}^{*}$ & $\underline{7.932}^{*}$ & $\underline{0.353}^{*}$ \\
独立训练 & $0.738^{*}$ & $\underline{0.335}^{*}$ & $3.259^{*}$ & $\underline{0.550}$ \\
集中式多任务 & $0.743^{*}$ & 0.358 & $2.662^{*}$ & $\underline{0.506}$ \\
直接顺序训练 & $\underline{0.348}^{*}$ & $\underline{0.270}$ & $\underline{6.922}^{*}$ & 0.568 \\
EWC & $\underline{0.659}^{*}$ & $\best{0.364^{*}}$ & 3.836 & $\underline{0.506}$ \\
SI & $\underline{0.461}^{*}$ & $\underline{0.303}^{*}$ & $\underline{8.186}^{*}$ & $\underline{0.555}^{*}$ \\
GPM & $\underline{0.458}^{*}$ & $\underline{0.290}^{*}$ & $\underline{7.095}^{*}$ & $\underline{0.556}^{*}$ \\
MER & \best{0.698} & $\underline{0.331}^{*}$ & $\best{3.611^{*}}$ & $\underline{0.532}$ \\
SAMCL & 0.697 & 0.352 & 3.855 & \best{0.563} \\
\bottomrule
\end{tabular}

\vspace{8pt}
\begin{tabular}{@{}lcccc@{}}
\toprule
\multicolumn{5}{c}{（b）后向变化 $R_{T,j}-R_{j,j}$} \\
方法 & OASIS Dice & 腹部 CT Dice & NLST TRE & 腹部 MR--CT Dice \\
\midrule
直接顺序训练 & $\underline{-0.395}^{*}$ & $\underline{-0.129}^{*}$ & $\underline{3.508}^{*}$ & 0  \\
EWC & $\underline{-0.083}^{*}$ & $\best{-0.011^{*}}$ & $\underline{0.386}$ & 0 \\
SI & $\underline{-0.281}^{*}$ & $\underline{-0.084}^{*}$ & $\underline{4.685}^{*}$ & 0 \\
GPM & $\underline{-0.281}^{*}$ & $\underline{-0.085}^{*}$ & $\underline{3.574}^{*}$ & 0 \\
MER & $\best{-0.036}$ & $\underline{-0.040}$ & $\best{0.077^{*}}$ & 0 \\
SAMCL & $-0.037$ & $-0.039$ & 0.306 & 0 \\
\bottomrule
\end{tabular}

\vspace{3pt}
\begin{minipage}{0.98\textwidth}
\thesistablefont
注：Dice 越高越好，NLST TRE（mm）越低越好。红色粗体标出 EWC、SI、GPM、MER 和 SAMCL 中的列最优值；下划线表示该值劣于 SAMCL，星号表示与 SAMCL 的配对 $t$ 检验达到 $p<0.05$。未配准、独立训练和集中式多任务不定义后向变化；最后任务的零后向变化不作最优标记。
\end{minipage}
\end{table}

从最终配准精度看，SAMCL 在 OASIS、腹部 CT--CT 和腹部 MR--CT 上分别取得 $0.697$、$0.352$ 和 $0.563$ 的 Dice，在 NLST 上取得 $3.855$ mm 的 TRE。相较 SI 和 GPM，SAMCL 在 OASIS 与腹部 CT--CT 上具有更高 Dice，并明显减小早期任务的后向下降。与 EWC 相比，SAMCL 在 OASIS 上具有更高最终 Dice 和更小后向下降；腹部 CT--CT 上 EWC 的最终 Dice 略高，且后向下降更小。该结果说明参数正则化在某一任务上仍可能有效，但其表现会随器官和损失形式变化。

在相同历史访问条件下，MER 与 SAMCL 均采用有限回放，因此二者的比较更直接。两种方法在 OASIS 最终 Dice 和前两个任务 BWT 上接近；SAMCL 在腹部 CT--CT 与第四个到达的腹部 MR--CT 上取得更高 Dice，MER 则在 NLST 上取得更低 TRE 和更小 TRE-BWT。由此可见，锐度感知元持续学习并未在所有任务上统一超过 MER，而是在不同任务之间形成了不同的稳定性--可塑性取舍。

独立训练和集中式多任务训练可以分别访问单一任务完整数据或全部任务数据。SAMCL 在 OASIS 上低于两项参考，在腹部 CT--CT 上接近集中式多任务训练，并在腹部 MR--CT 上高于独立训练和集中式多任务训练。直接顺序训练在最后任务上取得 $0.568$ 的 Dice，略高于 SAMCL 的 $0.563$；由于最后一个任务之后没有进一步训练，该差异反映当前任务适应，而不能用于评价历史保持。在前三个历史任务上，SAMCL 缩小了直接顺序训练与完整数据访问参考之间的差距；有限历史回放仍不能替代集中式访问全部历史任务。

图\ref{fig:samcl-qualitative}给出 OASIS 与腹部 CT--CT 的代表性图像对。每个任务的上排显示固定图像与配准后移动图像之间的差异图，下排显示不同方法得到的配准结果及对应 Dice。直接顺序训练、SI 和 GPM 在早期任务上出现较明显差异，而 SAMCL 的差异图和变形后图像更接近独立训练与集中式多任务训练。该可视化与表\ref{tab:samcl-main-results}中的早期任务保持结果一致。

\IfFileExists{figures/ch04/samcl_qualitative.pdf}{%
\begin{figure}[htbp]
\centering
\includegraphics[width=\textwidth]{ch04/samcl_qualitative.pdf}
\caption[不同训练策略的代表性配准结果]{不同训练策略在 OASIS 与腹部 CT--CT 代表性图像对上的配准结果。每个任务左侧给出固定图像与移动图像；各方法的上排为固定图像与配准后移动图像之间的差异图，下排为配准结果及对应 Dice。CT 图像显示时将强度裁剪到 $[-200,300]$。}
\label{fig:samcl-qualitative}
\end{figure}
}{}

\subsubsection{记忆容量消融与前向泛化分析}
\label{subsubsec:samcl-ablation-generalization}

图\ref{fig:samcl-ablation-generalization}(a)比较了 SAMCL 在经验回放容量分别为 $50$、$100$ 和 $200$ 时的配准结果。对于 OASIS、腹部 CT--CT 和腹部 MR--CT 三项以 Dice 评价的任务，增大回放容量总体上提高了最终配准精度。更大的回放集合能够保存更多历史图像对，使早期任务在后续阶段被重新采样的覆盖范围增加，因此有利于历史知识保持。该结果也表明，SAMCL 的性能依赖可保存的历史原始图像数量，属于有限回放条件下的结果。

除回放容量外，图\ref{fig:samcl-ablation-generalization}(b)考察锐度感知最小化对域内与域外前向泛化的影响。模型只在较早到达的任务上完成训练，随后直接在未参与训练的数据上测试。域内设置包括 CTCT$\rightarrow$CTCT 和 MRCT$\rightarrow$MRCT；域外设置包括 CTCT$\rightarrow$MRCT、CTCT$\rightarrow$NLST 和 OASIS$\rightarrow$NLST。加入 SAM 后，前三个以 Dice 评价的方向均取得更高结果，后两个以 TRE 评价的方向均取得更低误差。该结果支持锐度感知更新能够在这些预设方向上改善前向泛化，并说明平坦参数邻域不仅影响已见任务保持，也会影响模型面对后续分布时的直接适用性。

\begin{figure}[htbp]
\centering
\includegraphics[width=0.90\textwidth]{ch04/plot_memory_nnn.pdf}\\[-0.2em]
{\thesisfigurefont (a) 记忆容量消融实验}\par\vspace{5pt}
\includegraphics[width=0.90\textwidth]{ch04/plot_sam_new.pdf}\\[-0.2em]
{\thesisfigurefont (b) 域内与域外前向泛化实验}
\caption[SAMCL 的消融与前向泛化实验]{SAMCL 的消融与前向泛化实验。（a）采用不同回放容量时的配准结果，其中 OASIS、CTCT 和 MRCT 报告 Dice（越高越好），NLST 报告 TRE（越低越好）；（b）采用或不采用锐度感知最小化时的域内与域外前向泛化结果。CTCT 表示腹部 CT--CT 配准，MRCT 表示腹部 MR--CT 配准。}
\label{fig:samcl-ablation-generalization}
\end{figure}

本节实验采用固定任务顺序和单一主要骨干，没有系统交换四项任务的到达次序。结果表明，SAMCL 在既定顺序下能够缓解遗忘并改善多项前向泛化结果，但任务顺序仍可能影响最终配准表现。

\section{分割与配准的联系及适用边界}
\label{sec:task-cl-synthesis}

两项研究的共同点在于，历史原始样本只能有限保留，而当前任务仍会持续改变共享参数。第三章通过性能矩阵区分当前学习、历史保持和前向泛化；本章在分割中直接沿用这一矩阵，在配准中保留相同的阶段记录方式，但按任务分别使用 Dice 或 TRE。因此，共同的是问题组织和比较原则，具体预测对象、监督形式与指标仍由任务决定。

两种方法对历史样本的利用有所侧重。ZSDERpp 需要在稀疏标注下维护类别判别与空间结构，因此回放条目包含图像、涂鸦和参考特征，并重新计算弱监督目标。SAMCL 的目标是保持空间匹配能力，回放条目是历史图像对及相应任务监督，元更新和锐度感知优化用于协调任务之间的梯度。前者的关键设计落在监督与表征保持，后者落在更新过程与损失邻域；同属有限回放不代表两种设计可以未经验证地互换。

现有结果支持的是各自实验条件下的局部结论。ZSDERpp 在六中心 Domain-CL 中取得更好的最终平均 Dice 和历史保持，但其前向泛化不优于所有比较方法；SAMCL 缓解了直接顺序配准的早期任务退化，但与 MER 的比较随任务和指标而改变。分割使用二维切片、配准使用三维图像对，二者的回放容量、存储字节量和训练代价也不相同。因此，本章不把它们合成为一个总体性能分数，也不据此宣称已经建立统一的分割--配准算法。

在医学影像处理流程中，分割可以为配准提供结构对应约束，配准可以帮助对齐纵向影像和传播结构信息。这种任务互补性解释了进一步研究共享表征和联合更新的价值，但本章没有实现上述双向连接。与本章保留少量历史样本不同，第\ref{chap:fedsubmerge}章同时收紧历史和跨中心访问条件，研究如何以梯度子空间替代原始样本回放，并将知识保护扩展到多个客户端。

\section{本章小结}
\label{sec:task-cl-summary}

本章分别研究了涂鸦监督持续分割与持续医学图像配准。ScribbleCL 直接继承第三章的三类场景、数据、任务序列和评价指标，并将当前训练监督改为涂鸦。在此基础上，ZSDERpp 把历史特征保持和容量受限回放纳入同一目标，并在 Class-CL 中处理背景语义变化。已完成的 Domain-CL 实验表明，该方法改善了最终多域性能与旧域保持。SAMCL 结合有限图像对回放、一阶元更新与锐度感知优化，在四任务序列中缓解了直接顺序训练的早期任务性能下降，并改善了部分前向泛化结果。
